% Ad Familiares 9.24 (ad-familiares-09-24) --- English translation
% Translated by Claude (Anthropic) from the Latin (Perseus).
% Style guide: STYLE.md.

\ciceroLetterOpener{CICERO PAETO S. D.}

\ciceroSection{1} That Rufus of yours, your friend, about whom you now write to me a second time --- I would help him as far as I could even if he had injured me, seeing how strenuously you exert yourself on his behalf; but in fact, since I gather both from your letters and from those he himself has sent me, and I judge accordingly, that my safety has been a matter of great concern to him, I cannot but be his friend --- not only on your recommendation, which carries the very greatest weight with me, as it should, but also out of my own will and judgment. For I want you to know, my Paetus, that the beginning of my suspicion and caution and care was your letter, with which other letters of many people afterwards agreed. For at Aquinum and at Fabrateria plots were laid against me --- you, I see, have heard about them --- and, as though they had divined how much trouble I was going to be to them, they did nothing other than try to crush me. Unsuspecting all this, I would have been less on my guard had I not been forewarned by you. So that friend of yours needs no recommendation with me. May the fortune of the Republic be such that he can come to know me as supremely grateful! But enough of this.

\ciceroSection{2} That you have stopped going out to dinners distresses me; for you have deprived yourself of a great delight and pleasure. Then too I am afraid --- if I may speak the truth --- that you may unlearn and forget some skill you used to have, of putting on little dinners yourself. For if even when you had models to imitate you were not making much progress, what am I to suppose you will do now? Spurinna at any rate, when I had laid the case before him and set out your former way of life, was indicating great peril to the highest interest of the Republic unless you should return to your earlier habits at the time the Favonius blew; at this season the thing might be borne, supposing that you cannot bear the cold.

\ciceroSection{3} But, by Hercules, my Paetus, jokes aside, I urge you on a matter that I think bears on the happy life: live with good, agreeable men, who love you. Nothing is more fitting for life, nothing better suited to living happily. And I refer it not to pleasure but to the common share of life and sustenance, and to the unbending of the spirit, which is most produced by intimate conversation --- and that is at its sweetest at table. Our forefathers were wiser in this than the Greeks: they speak of \textit{symposia} or \textit{syndeipna} \textit{[Greek: symposia ... syndeipna]} --- that is, ``drinkings-together'' or ``dinings-together'' --- but we speak of \textit{convivia}, ``livings-together,'' because then most of all is life shared. You see how, in philosophizing, I try to call you back to dinners. Take care of yourself; that you will most easily accomplish by dining out.

\ciceroSection{4} But beware, if you love me, of supposing that, because I write somewhat jocosely, I have flung off care for the Republic. Be persuaded of this, my Paetus: day and night I do nothing else, care for nothing else, except that my fellow citizens should be safe and free. No place do I leave untried to warn, to act, to foresee. In short, I am of this mind: if in this care and administration my life must be laid down, I shall reckon that things have turned out splendidly for me. Again and again, farewell.
